\documentclass[10pt]{article}
\usepackage[a4paper,margin=22mm]{geometry}
\usepackage{amsmath,amssymb}
\usepackage[hidelinks]{hyperref}
\usepackage{xcolor}

\newcommand{\DN}{\underline{DN}}
\newcommand{\Om}{\Omega}

\begin{document}
\begin{center}
{\LARGE\bfseries The Finite-Type Diametral-Dimension Question\\[2pt]
Has a Known Negative Answer\par}
\vspace{8pt}
{\large Literature-status packet for arXiv:1908.01838\par}
\vspace{3pt}
27 August 2026
\end{center}
\vspace{3pt}

\begin{center}
\fcolorbox{black}{gray!10}{\parbox{0.86\textwidth}{
\textbf{Status: literature already answered.}
The later paper arXiv:2210.13593 explicitly answers the finite-type part of
Question 1.1 of arXiv:1908.01838 by a counterexample.
}}
\end{center}

\section{Original question}

In \emph{Some Remarks on Diametral Dimension and Approximate Diametral
Dimension of Certain Nuclear Fr\'echet Spaces}~\cite{source}, Do\u{g}an asks
in Question 1.1 (PDF page 2): for a nuclear Fr\'echet space $E$ satisfying
$\DN$ and $\Om$, if its diametral dimension $\Delta(E)$ coincides with that
of a power series space, must its approximate diametral dimension $\delta(E)$
also coincide, and conversely?

The source proves a full affirmative equivalence for infinite-type power
series spaces. For finite type it proves
\[
 \delta(E)=\delta(\Lambda_1(\varepsilon))
 \quad\Longleftrightarrow\quad
 \bigl[\Delta(E)=\Delta(\Lambda_1(\varepsilon))
 \text{ and $E$ has a prominent bounded set}\bigr]
\]
(Theorem 4.8, PDF page 11). Thus the remaining issue is whether equality of
the finite-type diametral dimensions alone forces the equality of approximate
diametral dimensions, or equivalently forces a prominent bounded set.

\section{Explicit later answer}

The later paper \emph{On Power Series Subspaces of Certain Nuclear Fr\'echet
Spaces}~\cite{answer} repeats the earlier problem as Question 1.2 and states
on PDF page 2 that the finite-type answer is negative. For every unstable
finitely nuclear exponent sequence $\alpha$, it constructs a nuclear
Fr\'echet K\"othe space $K_\alpha$ satisfying $\DN$ and $\Om$.

Propositions 3.6 and 3.7 (PDF pages 20--22) prove, with
$\varepsilon_n=\alpha_{n+1}$,
\[
 \Delta(K_\alpha)=\Delta(\Lambda_1(\varepsilon)),
 \qquad
 \delta(K_\alpha)\ne\delta(\Lambda_1(\varepsilon)).
\]
Remark 3.8 explicitly identifies these two propositions as a negative answer
to Question 1.2. Theorem 4.8 (PDF page 26) records the equivalent obstruction:
$K_\alpha$ has no prominent bounded set. The answering paper cites the
original paper and clearly identifies the question it resolves.

\section{Classification and scope}

This is a separate later paper explicitly answering the source question, so
the correct run classification is \texttt{literature\_already\_answered}, not
a new counterexample. The original infinite-type equivalence and the
finite-type reverse implication remain intact. The counterexample refutes
only the unrestricted finite-type implication
\[
 \Delta(E)=\Delta(\Lambda_1(\varepsilon))
 \ \Longrightarrow\ 
 \delta(E)=\delta(\Lambda_1(\varepsilon)).
\]

\begin{thebibliography}{9}
\raggedright
\bibitem{source}
N.~Do\u{g}an,
\emph{Some Remarks on Diametral Dimension and Approximate Diametral
Dimension of Certain Nuclear Fr\'echet Spaces},
Bull. Belg. Math. Soc. Simon Stevin \textbf{27} (2020), 353--368;
arXiv:1908.01838v2.

\bibitem{answer}
N.~Do\u{g}an,
\emph{On Power Series Subspaces of Certain Nuclear Fr\'echet Spaces},
arXiv:2210.13593v1 (2022).
\end{thebibliography}

\end{document}
